% Word count aim                0800-0900
% Current: 997
\chapter{Supplementary Software Design}
\label{cha:design}
This chapter discusses the supplementary software developed for
carrying out multiple experiments on \textit{simulated} \acrshort{bbo}
problems efficiently. Trivially, analysis cannot be carried out on
truly black-box instances; the problems are emulated to be black-box
by providing only the \textit{zeroth-order} information.

\section{Purpose of the Software}
% Max 200 words
% Explain why I made this software. Mention the motivations for making it.
A supplementary experimental framework, \texttt{regret}, was developed to
support the orchestration and analysis of experiments on pseudo-Boolean and
combinatorial optimization problems. The framework was motivated by the
absence of readily available tooling that combines (i)~flexible, declarative
batch experiment execution, (ii)~trajectory-level data collection, and
(iii)~regret-specific analysis pipelines in a single cohesive system.

Existing platforms such as Google Vizier~\cite{google_vizier} and
EvoX~\cite{evox} address hyperparameter optimization and population-based
search respectively, but neither provides native support for regret metrics,
runtime profile analysis, or the configuration-first experiment
execution model required
for this project's experimental protocol. \texttt{regret} was therefore
purpose-built to fill this gap, providing a reproducible, extensible
environment for studying simple, instantaneous, and \acrshort{cr} across
a controlled suite of algorithms and benchmark problems.

\subsection{Preexisting software}
% Further mention Google Vizier, Evox and other commericially used software
Several established frameworks were considered prior to development.
\textbf{Google Vizier}~\cite{google_vizier} provides a scalable black-box
optimization service with a client--server architecture, but is oriented
toward hyperparameter tuning rather than controlled algorithmic benchmarking,
and does not expose trajectory-level data in a form amenable to regret
analysis. \textbf{EvoX}~\cite{evox} is a JAX-based
evolutionary computation
library offering GPU-accelerated population methods, but its abstraction layer
couples algorithm state to framework internals in ways that preclude the
lightweight callback-driven trajectory recording required here.
General-purpose scientific workflow systems such as
Snakemake~\cite{snakemake} and Luigi~\cite{luigi_software} offer
pipeline orchestration but impose significant configuration overhead for the
combinatorial (algorithm~$\times$~problem~$\times$~budget) sweeps central to
this study.

\subsection{Capabilities}
% Max 250 words
% Describe how this software can be utilized to conduct experiments
% Highlight the dispatching capability with the help of YAML configurations
% This can help with flexible dispatch, parallelized workflow and
% batching of experiment orchestrations, all from the flexibility of a config.
The framework supports the full experimental life-cycle through a single
\acrfull{cli}, \texttt{run\_experiment}, which accepts one or more
\acrshort{yaml} configuration files and executes the appropriate sub-command.
Table~\ref{tab:cli-commands} summarizes the available commands.

\small
\begin{longtable}{p{3.8cm}p{9.2cm}}
  \hline
  \textbf{Command} & \textbf{Description} \\
  \hline
  \texttt{validate} &
  Parse and validate one or more \acrshort{yaml} configurations
  against the \acrshort{json} Schema and
  semantic registry checks without executing any computation. \\[4pt]
  \texttt{plan} &
  Print a dry-run summary: problem/algorithm combinations, total runs, and
  evaluation budgets. \\[4pt]
  \texttt{run [\texttt{-{}-no-plot}]} &
  Execute one or more experiment configurations; optionally suppress plot
  generation after completion. \\[4pt]
  \texttt{analyze} &
  Regenerate all plots from previously saved \acrshort{json} results without
  re-running experiments. \\[4pt]
  \texttt{table [\texttt{-{}-format}]} &
  Generate summary tables from existing results in \texttt{latex},
  \texttt{csv}, or \texttt{markdown} format. \\
  \hline
  \caption[CLI sub-commands]{\acrshort{cli} sub-commands exposed by the
    \texttt{run\_experiment} entry
  point (equivalently \texttt{python -m regret}).}
  \label{tab:cli-commands}
\end{longtable}
% \end{table}

\paragraph{Batch Experiment Execution.}
All experiments are specified declaratively via \acrshort{yaml}
configuration files
validated against a \acrshort{json} Schema prior to execution. A
configuration defines
the suite name, evaluation budgets, run count, output mode, and lists of
problem and algorithm specifications. Problems and algorithms are referenced
by string keys resolved against internal registries, providing a clean
separation between configuration and implementation.

\paragraph{Parallelized execution.}
When \texttt{parallel: true} is set in the \acrshort{yaml} configuration,
independent runs are initiated via
Python's \texttt{ProcessPoolExecutor}, with per-run random seeds fixed to
indices $0, 1, \ldots, r-1$ to ensure full reproducibility. Results are
reordered by seed index after completion.

\paragraph{Flexible output modes.}
In \texttt{lite} mode, only summary statistics are persisted per run. In
\texttt{full} mode, the complete evaluation trajectory (a sequence of
\texttt{(evaluation\_index, current\_value, best\_value)} tuples) is
recorded and stored alongside aggregate statistics in structured
\acrshort{json} files.

\paragraph{Batched multi-configuration execution.}
All sub-commands accept multiple configuration files, enabling sequential
execution of logically grouped experiment suites from a single invocation.

\subsection{Design Decisions}
% Max 150 words
% Why the dispatching framework? Why the
The software was intended to be extensive and configurable for other
developers and users to interact with quickly, with options for
modularity when required.
This was enabled by the adoption of configuration-first, job dispatch
policy for independent and declarative experimentation. Problems and
Algorithms to test against can easily be defined by extending
abstract classes \texttt{Problem} and \texttt{Algorithm}, ensuring a
common interface for consistent interactions.

% \begin{figure}[h]
% \centering
% \begin{minipage}{0.48\textwidth}
% \textbf{Traditional}
% \begin{verbatim}
% {
%   ...
%   "results": [
%     {
%       "regret": --,
%       "best_value: --,
%       "optimal": --,
%       "evaluations": --,
%       "seed": --,
%       "trajectory": [
%         [
%           1,
%           618324347019193.0,
%           618324347019193.0
%         ],
%         [
%           2,
%           622722393530297.0,
%           622722393530297.0
%         ],
%         [
%           3,
%           622722393538489.0,
%           622722393538489.0
%         ],
%         ...
%         ]
%     ...
%     ]
% ...
% }
% \end{verbatim}
% \end{minipage}
% % \hfill
% \begin{minipage}{0.48\textwidth}
% \textbf{Version B}
% \begin{verbatim}
% def add_numbers(x, y):
%     result = x + y
%     return result
% \end{verbatim}
% \end{minipage}
% \end{figure}

\subsubsection{Independent Job Execution framework}

The \acrshort{cli} is implemented using \textbf{Typer}~\cite{typer}, which maps
sub-commands to typed Python functions with automatic argument parsing. The
``orchestration'' layer coordinates
problem and algorithm instantiation, result persistence, and plot
generation, delegating individual runs to \texttt{ExperimentRunner}. Parallel
execution uses \texttt{ProcessPoolExecutor} from Python's standard
\texttt{concurrent.futures} module; results are collected via
\texttt{as\_completed} and reordered by seed for deterministic output
ordering. Figure~\ref{fig:dispatch-pipeline} illustrates the full validation
and the experiment (job) dispatch pipeline.

% TODO: Expand flowchart
\begin{figure}[H]
  \centering
  \begin{tikzpicture}[
      node distance = 2cm and 3cm,
      on grid,
      box/.style = {
        rectangle, rounded corners=3pt, draw=black!70,
        fill=blue!8, minimum width=1.8cm,
        minimum height=0.8cm, font=\small, align=center
      },
      decision/.style = {
        diamond, draw=black!70, fill=orange!15,
        minimum width=2.2cm, minimum height=1.0cm,
        font=\scriptsize, align=center, aspect=2
      },
      errbox/.style = {
        rectangle, draw=red!60, fill=red!5, rounded corners=2pt,
        font=\scriptsize, align=center
      },
      arrow/.style = {->, >=Stealth, thick},
      darrow/.style = {->, >=Stealth, thick, dashed, red!60}
    ]

    \node[box]   (yaml)  {\acrshort{yaml}\\configuration};
    \node[box,   right=of yaml]  (schema) {Schema\\Validation};
    \node[box,   below=of schema] (sem)  {Semantic\\Validation};
    \node[box,   left=of sem]  (parse) {Parse Problems\\and Algorithms};

    \node[decision, below=of parse] (par) {parallel?};

    \node[box, right=3.0cm of par] (seq) {Sequential\\Runner};
    \node[box, left=3.0cm of par] (pool) {ProcessPool\\Executor};

    \node[box, below=of par] (persist) {Persist \acrshort{json}\\Results};
    \node[box, below=of persist] (plots) {Generate\\Plots};

    \node[errbox, right=of schema] (e1) {\textit{ValidationError}};
    \node[errbox, right=of sem]  (e2) {\textit{ValidationError}};

    \draw[arrow] (yaml)  -- (schema);
    \draw[arrow] (schema) -- (sem);
    \draw[arrow] (sem)  -- (parse);
    \draw[arrow] (parse) -- (par);

    \draw[arrow] (par) -- node[above left] {\scriptsize no} (seq);
    \draw[arrow] (par) -- node[above right] {\scriptsize yes} (pool);

    \draw[arrow] (seq) |- (persist);
    \draw[arrow] (pool) |- (persist);

    \draw[arrow] (persist) -- (plots);

    \draw[darrow] (schema) -- (e1);
    \draw[darrow] (sem)  -- (e2);

  \end{tikzpicture}
  \caption[Configuration parse pipeline]{Validation and job dispatch
    pipeline. \acrshort{yaml} configuration files are
    validated structurally via \acrshort{json} Schema and then
    semantically against
    class registries before execution begins. Any validation failure
    raises a \texttt{ValidationError} and halts the pipeline before
  computation begins.}
  \label{fig:dispatch-pipeline}
\end{figure}

\subsubsection{Raw Output \acrshort{json} files}
The output \acrshort{json} files are \textit{compacted} by replacing
redundant stylistic whitespaces and newlines with other separators.
This reduced the raw file sizes by up to 50\%.
\acrshort{json} files were preferred for storing the raw files since
it included \textit{mixed} data, including metadata, general
statistics for each run, along with the trajectories.
A `NoSQL' (key-value) database was considered but not used since the
experiments were relatively small scale.

\subsubsection{UML Diagram}
% Diagram
Figure~\ref{fig:uml} presents the core class structure of the
framework. Abstract base classes \texttt{Problem} and
\texttt{Algorithm} define the shared interface enforced across all concrete
implementations. Concrete subclasses are registered by string key and
resolved at runtime. \texttt{ExperimentRunner} orchestrates individual runs
and delegates result persistence, while the \texttt{CoolingSchedule} protocol
is resolved from \acrshort{yaml} at parse time via a \textit{registry}.

\begin{figure}[H]
  \centering
  \includegraphics[width=\linewidth]{uml.png}
  \caption[UML class diagram]{Simplified UML class diagram of the
    \texttt{regret} framework.
    Solid triangle arrows ($\triangleright$) denote inheritance;
    dashed arrows denote dependency or instantiation.
    Hashes `\#' indicate abstract declarations.
    Several concrete subclasses (e.g.\ \texttt{MuPlusLambdaEA},
      \texttt{\acrlong{rlse}}, \texttt{NKLandscape},
  \texttt{LinearCooling}) are omitted for clarity.}
  \label{fig:uml}
\end{figure}

\subsection{Limitations}
% Does not support multi-objective optimization, continuous
% optimization, limited to binary strings.
The framework is currently constrained to \textbf{binary search spaces}
$\{0,1\}^n$; continuous and non-binary integer decision domains are not
supported.
\textbf{Multi-objective optimization} is not implemented. The NK landscape
optimum computation uses exhaustive search, restricting its practical use to
$n \leq 20$. Runtime profile analysis assumes integer-valued or explicitly
discretized fitness levels for the exact tail-sum identity:
\[
  \mathbb{E}[\mathrm{CR}(T)]
  \;=\;
  \sum_{v=1}^{f^{*}}\sum_{t'=1}^{T}
  \bigl[1 - P(\tau_v \leq t')\bigr];
\]
real-valued fitness functions (e.g.\ \texttt{NK-Landscape}) require
discretization and therefore yield approximate runtime-profile results.

% Considering addition to the next chapter.
\subsection{Result Interpretations}
\label{subsec:resint}
% % Max 200 words
% % Describe the output file structure, explaining the organisation.

All experiment output is written under a structured directory hierarchy.
Raw \acrshort{json} result files are stored at:

\texttt{results/raw/<suite>/<problem>/<algorithm>/n<N>/b<B>.json}

\noindent where \texttt{<suite>}, \texttt{<problem>}, and
\texttt{<algorithm>} are
normalized display names, \texttt{N}~is the problem dimension, and \texttt{B}
is the evaluation budget. Each file contains three top-level keys:
\texttt{metadata} (algorithm, problem class, dimension, budget, run count,
mode, and timestamp), \texttt{statistics} (mean, median, standard deviation,
quartiles, and probability of optimum across all runs), and \texttt{results}
(an ordered list of per-run dictionaries each containing regret, best value,
  optimality flag, evaluation count, seed, and the complete trajectory
(in \texttt{full} mode)).

Generated figures are written under \texttt{results/figures/} following an
analogous hierarchy:

\texttt{results/figures/<suite>/<problem>/n<N>/<subdir>/<filename>.pdf}

\noindent Plot subdirectories are configurable via
\texttt{layout.structure} in the YAML
configuration; the defaults are \texttt{aggregate/} for cross-budget plots,
\texttt{history/} for trajectory and regret time-series plots, and
\texttt{distributions/} for budget-specific box-plots and performance profiles.
Runtime profile plots are written to a \texttt{profiles/} subdirectory within
the raw results tree. All figures are saved as \acrshort{pdf} at
300\,\acrshort{dpi}.
Table~\ref{tab:output-structure} summarizes the full output organization.

\small
\begin{longtable}{p{0.55\linewidth}p{0.45\linewidth}}
  \hline
  \textbf{Path pattern} & \textbf{Contents} \\
  \hline
  \multicolumn{2}{l}{\textit{Raw data (JSON)}} \\
  \hline
  \small\texttt{results/raw/\newline<suite>/<prob>/<algo>/n<N>/b<B>.json}
  & Per-budget result file: metadata, aggregate statistics, per-run regrets,
  and (in \texttt{full} mode) full evaluation trajectories. \\[4pt]
  \small\texttt{results/raw/<suite>/<prob>/}\newline\texttt{profiles/n<N>/cr\_profile\_data\_<algo>.csv}
  & Runtime profile CSV; columns: \texttt{evaluation},
  \texttt{expected\_cumulative\_regret},
  \texttt{mean\_cumulative\_regret}. \\
  \hline
  \multicolumn{2}{l}{\textit{Figures (\acrshort{pdf}, 300\,DPI)}} \\
  \hline
  \texttt{.../n<N>/aggregate/}
  & Cross-budget plots: regret curves, convergence probability,
  comparison heatmap. \\[4pt]
  \texttt{.../n<N>/history/}
  & Trajectory plots: best/current value, instantaneous and cumulative
  regret, \acrfull{ttfo} distribution. \\[4pt]
  \texttt{.../n<N>/distributions/}
  & Budget-specific plots: regret box-plots and performance profiles. \\[4pt]
  \texttt{.../n<N>/profiles/}
  & Runtime profile plots: inverse profile surface/curves,
  $\mathbb{E}[\mathrm{CR}(T)]$ verification.
  Requires \texttt{suite.profile: true}. \\
  \hline
  \caption[Experiment output organization]{Output file organization
    for a representative experiment.
    ``Normalization'' replaces spaces, parentheses, and special characters
  with underscores or removes them.}
  \label{tab:output-structure}
\end{longtable}

\small
\begin{longtable}{p{4cm}p{8cm}p{2cm}}
  \hline
  \textbf{Plot} & \textbf{Output filename pattern - \texttt{\acrshort{pdf}}
  format} & \textbf{Profile required} \\
  \hline
  Regret curves            & \texttt{regret\_curves}
  & No \\
  Convergence probability  & \texttt{convergence\_probability}
  & No \\
  Comparison heatmap       & \texttt{comparison\_heatmap}
  & No \\
  Regret box-plots          & \texttt{regret\_boxplot\_b\{budget\}}
  & No \\
  Performance profile      &
  \texttt{performance\_profile\_b\{budget\}}    & No \\
  History (best-so-far)    & \texttt{history\_best}
  & No \\
  Instantaneous regret     &
  \texttt{history\_regret\_instantaneous}       & No \\
  Cumulative regret        & \texttt{history\_regret\_cumulative}
  & No \\
  \acrshort{ttfo} distribution        &
  \texttt{history\_ttfo\_markers}               & No \\
  Runtime profile surface  &
  \texttt{inverse\_runtime\_profile\_surface\_\{alg\}} & Yes \\
  Runtime profile curves   &
  \texttt{inverse\_runtime\_profile\_curves}   & Yes \\
  CR profile verification  & \texttt{cr\_profile\_verification}
  & Yes \\
  \hline
  \caption[Generated experiment plot types]{Plot types generated
    across experiments. ``Profile required'' denotes
  plots that require \texttt{suite.profile:~true}}
  \label{tab:plot-types}
\end{longtable}

\section{Testing strategy}
% Mention TDD, Red-Green Testing, selective coverage (weighted)
\begin{wrapfigure}{r}{0.5\textwidth} %this figure will be at the right
  \centering
  \includegraphics[width=0.5\textwidth]{red-green-refactor.png}
  \caption[Red--Green--Refactor \acrshort{tdd}
  technique]{Red--Green--Refactor \acrshort{tdd} technique. Start by
    writing a failing test (the red step). Next, write code that is
    just enough to make the written test pass (the green step).
    Finally, as the regularities and properties of the problem space
    are understood, refactors can be implemented that keeps the code
    clean (the blue step). The refactor step is safe since it will
  always be backed up by tests.~\cite{calais2023test}}
\end{wrapfigure}
The test suite follows a \acrfull{tdd} discipline
using a red--green--refactor workflow: tests were written prior to or alongside
implementation and serve as executable
specifications~\cite{agha2023test,calais2023test}. Coverage is measured
using the \texttt{coverage}~\cite{coverage} package with branch
tracking enabled, and a
minimum threshold of \textbf{50\% branch coverage} is enforced in
\acrfull{ci} via GitHub Actions. Coverage is also computed across
Python~3.12, 3.13, and 3.14 in \acrfull{ci} for redundancy and
support (the project must be compatible for \texttt{python} $\geq$ 3.12).

Test organization follows a selective weighting strategy: core algorithmic
components (\texttt{test\_algorithms.py}, \texttt{test\_metrics.py},
\texttt{test\_problems.py}) and integration paths
(\texttt{test\_integration.py}) receive dense coverage, while infrastructure
components such as CLI wiring and plot dispatch are covered at a coarser
granularity.

\section{Project Dependencies}
The project is built with \texttt{hatchling}~\cite{hatchling} and managed
with \texttt{uv}~\cite{uv}, a Rust-based package and environment manager that
replaces the Python standard \texttt{pip}/\texttt{venv}.
Runtime installation: \texttt{uv sync}; development installation:
\texttt{uv sync -{}-group dev}.

\small
\begin{longtable}{p{3cm}p{3cm}p{8cm}}
  \hline
  \textbf{Package} & \textbf{Min.\ version} & \textbf{Purpose} \\
  \hline
  \multicolumn{3}{l}{\textit{Scientific computing}} \\
  \hline
  \texttt{numpy}~\cite{numpy}      & 2.4.3  & Array operations,
  random number generation \\
  \texttt{scipy}~\cite{scipy}      & 1.16.3 & Statistics, drift
  computations, profile integrals \\
  \texttt{pandas}~\cite{pandas}    & 2.3.3  & Results aggregation,
  tabulation, CSV export \\
  \hline
  \multicolumn{3}{l}{\textit{Visualisation}} \\
  \hline
  \texttt{matplotlib}~\cite{matplotlib} & 3.10.8 & All static plot
  generation \\
  \texttt{ipympl}                       & 0.10.0 & Interactive
  Matplotlib in notebooks \\
  \texttt{networkx}~\cite{networkx}& 3.6.1  & Graph problem
  instances (e.g.\ Petersen colouring) \\
  \hline
  \multicolumn{3}{l}{\textit{Configuration and \acrshort{cli}}} \\
  \hline
  \texttt{pyyaml}~\cite{pyyaml}      & 6.0.3  & Experiment
  configuration parsing \\
  \texttt{jsonschema}~\cite{jsonschema}  & 4.26.0 & \acrshort{yaml}
  schema validation \\
  \texttt{typer}~\cite{typer}       & 0.24.1 &
  \texttt{run\_experiment} \acrshort{cli} \\
  \texttt{tqdm}~\cite{tqdm}        & 4.67.1 & Progress bars for
  parallel runs \\
  \hline
  \multicolumn{3}{l}{\textit{Notebook and interactive demonstration}} \\
  \hline
  \texttt{jupyter}~\cite{jupyter}     & 1.1.1  & Notebook server \\
  \texttt{ipykernel}   & 7.2.0  & Jupyter kernel support \\
  \texttt{ipywidgets}  & 8.1.8  & Notebook UI widgets \\
  \hline
  \multicolumn{3}{l}{\textit{Development --- linting (optional)}} \\
  \hline
  \texttt{ruff}~\cite{ruff} & --- & Fast Rust-based linter and
  formatter (E, W, F, I, B, UP rules) \\
  \hline
  \multicolumn{3}{l}{\textit{Development --- testing (optional)}} \\
  \hline
  \texttt{pytest}~\cite{pytest}    & --- & Unit and integration test runner \\
  \texttt{coverage}~\cite{coverage}  & --- & Branch coverage
  reporting ($\geq$50\% enforced) \\
  \texttt{tabulate}  & --- & Tabular output in test fixtures \\
  \hline
  \multicolumn{3}{l}{\textit{Development --- type checking (optional)}} \\
  \hline
  \texttt{ty}~\cite{ty} & --- & Rust-based static type checker (Astral) \\
  \hline
  \multicolumn{3}{l}{\textit{Development --- documentation (optional)}} \\
  \hline
  \texttt{zensical}              & --- & Documentation site generator \\
  \texttt{mkdocstrings-python}   & --- & Auto-generated API
  reference from docstrings \\
  \hline
  \caption[Table of software prerequisites]{%
    Package dependencies for the \texttt{regret} framework (Python~$\geq$~3.12).
  }
  \label{tab:dependencies}
\end{longtable}
